\chapter*{Glossari de símbols}
\addcontentsline{toc}{chapter}{Glossari de símbols}
\label{cap:glossari-simbols}
\markboth{Glossari de símbols}{Glossari de símbols}

Aquest glossari recull la notació pròpia del llibre, agrupada per blocs
temàtics. La darrera columna indica el capítol on el símbol s'introdueix o es
defineix.

\begingroup
\renewcommand{\arraystretch}{1.35}

\subsection*{Llenguatge bàsic: aplicacions, dominis i composició}

\noindent
\begin{longtable}{@{}p{0.20\textwidth}p{0.62\textwidth}p{0.10\textwidth}@{}}
\toprule
\textbf{Símbol} & \textbf{Significat} & \textbf{Cap.}\\
\midrule
\endfirsthead
\toprule
\textbf{Símbol} & \textbf{Significat} & \textbf{Cap.}\\
\midrule
\endhead
\bottomrule
\endfoot
\(\Dom(f,a)\) & \(a\) és el domini de l'aplicació \(f\) & \ref{cap:aplicacions}\\
\(\Cod(f,b)\) & \(b\) és el codomini de l'aplicació \(f\) & \ref{cap:aplicacions}\\
\(\Com(f,g,h)\) & \(h\) és la composició de \(f\) i després \(g\) & \ref{cap:aplicacions}\\
\(g\circ f\) & composició d'aplicacions (primer \(f\), després \(g\)) & \ref{cap:aplicacions}\\
\(\id_A\), \(\mathrm{id}\) & aplicació identitat sobre \(A\) & \ref{cap:aplicacions}\\
\(f\colon A\to B\) & aplicació de domini \(A\) i codomini \(B\) & \ref{cap:aplicacions}\\
\(x\colon X\to A\) & element generalitzat de \(A\) a l'estadi \(X\) & \ref{cap:aplicacions}\\
\(\mathbf1\) & conjunt terminal (un sol punt) & \ref{cap:aplicacions}\\
\(\mathbf0\) & conjunt inicial (buit) & \ref{cap:aplicacions}\\
\(a\colon\mathbf1\to A\) & element global o punt de \(A\) & \ref{cap:aplicacions}\\
\(A\cong B\) & \(A\) i \(B\) són isomorfs & \ref{cap:aplicacions}\\
\end{longtable}

\subsection*{Parts, lògica interna i construccions universals}

\noindent
\begin{longtable}{@{}p{0.20\textwidth}p{0.62\textwidth}p{0.10\textwidth}@{}}
\toprule
\textbf{Símbol} & \textbf{Significat} & \textbf{Cap.}\\
\midrule
\endfirsthead
\toprule
\textbf{Símbol} & \textbf{Significat} & \textbf{Cap.}\\
\midrule
\endhead
\bottomrule
\endfoot
\(S\hookrightarrow A\) & monomorfisme; \(S\) és una part de \(A\) & \ref{cap:parts}\\
\(P(A)\) & conjunt de parts de \(A\) & \ref{cap:parts-interna}\\
\(A\times B\) & producte de conjunts & \ref{cap:parts}\\
\(A+B\) & suma (coproducte) de conjunts & \ref{cap:parts}\\
\(B^A\) & exponencial: aplicacions de \(A\) cap a \(B\) & \ref{cap:exponencials}\\
\(a\mathrel{\varepsilon}\beta\) & pertinença generalitzada interna & \ref{cap:parts-interna}\\
\(\neg_A\) & complement intern sobre \(P(A)\) & \ref{cap:parts-interna}\\
\(\exists_f,\ \forall_f\) & quantificadors interns al llarg de \(f\) & \ref{cap:parts-interna}\\
\(\abs{A}\) & cardinal (mida) de \(A\) & \ref{cap:cardinals}\\
\end{longtable}

\subsection*{Naturals, ordre i transfinit}

\noindent
\begin{longtable}{@{}p{0.20\textwidth}p{0.62\textwidth}p{0.10\textwidth}@{}}
\toprule
\textbf{Símbol} & \textbf{Significat} & \textbf{Cap.}\\
\midrule
\endfirsthead
\toprule
\textbf{Símbol} & \textbf{Significat} & \textbf{Cap.}\\
\midrule
\endhead
\bottomrule
\endfoot
\(\N,\ \Z,\ \Q,\ \R\) & naturals, enters, racionals, reals & \ref{cap:infinitud}\\
\(\seg{n}\) & segment inicial \(\{0,\dots,n-1\}\) de \(\N\) & \ref{cap:infinitud}\\
\(A\preceq B\) & \(A\) s'injecta en \(B\) (mida menor o igual) & \ref{cap:cantor}\\
\(A\prec B\) & \(A\) té mida estrictament menor que \(B\) & \ref{cap:cantor}\\
\((S,R)\) & bon ordre: suport \(S\) amb relació \(R\) & \ref{cap:zorn}\\
\(\WO(A)\) & bons ordres sobre parts de \(A\) & \ref{cap:zorn}\\
\(\camp(S,R)\) & suport d'un bon ordre (és \(S\)) & \ref{cap:zorn}\\
\(\Hart(A)\) & conjunt de Hartogs de \(A\) & \ref{cap:zorn}\\
\([(S,R)],\ [R]\) & tipus d'ordre d'un bon ordre & \ref{cap:zorn}\\
\(\omega\) & tipus d'ordre dels naturals & \ref{cap:ordinals}\\
\(\aleph_\alpha\) & cardinals infinits (alephs) & \ref{cap:cardinals}\\
\(\aleph_0\) & cardinal dels naturals; \(\abs{\N}\) & \ref{cap:cardinals}\\
\(\kappa^{+}\) & cardinal successor de \(\kappa\) & \ref{cap:cardinals}\\
\(\beth_\alpha\) & cardinals beth & \ref{cap:alephs}\\
\(\mathfrak c\) & el continu; \(\abs{P(\N)}=2^{\aleph_0}\) & \ref{cap:alephs}\\
\(\operatorname{cf}(\kappa)\) & cofinalitat de \(\kappa\) & \ref{cap:alephs}\\
\end{longtable}

\endgroup
